\begin{table}[H]
\centering
\caption{Productividad laboral departamental, 2009--2024pr}
\label{tab:dept_productividad_resumen_24}
\scriptsize
\begin{tabular}{lrrrr}
\toprule
Departamento & PIB/trab. 2024pr & Crec. & PIB/hora 2024pr & Crec. \\
\midrule
Chocó & 28,6 & 3,5\% & 11,9 & 4,0\% \\
Caquetá & 24,9 & 2,9\% & 10,5 & 3,8\% \\
Quindío & 33,8 & 2,2\% & 15,1 & 3,0\% \\
Risaralda & 38,7 & 2,3\% & 16,4 & 2,8\% \\
Cauca & 26,6 & 1,4\% & 13,7 & 1,8\% \\
Bolívar & 45,1 & 0,9\% & 19,8 & 1,8\% \\
Santander & 62,0 & 1,5\% & 26,1 & 1,8\% \\
Bogotá D.C. & 67,4 & 1,3\% & 28,6 & 1,7\% \\
Valle del Cauca & 54,5 & 1,4\% & 23,4 & 1,6\% \\
Tolima & 40,0 & 1,3\% & 16,9 & 1,6\% \\
Nariño & 17,8 & 0,4\% & 9,1 & 1,2\% \\
Meta & 63,1 & 0,8\% & 25,9 & 1,2\% \\
Caldas & 36,6 & 0,6\% & 15,6 & 1,2\% \\
Sucre & 22,0 & 0,2\% & 10,4 & 1,0\% \\
Antioquia & 47,7 & 0,9\% & 20,3 & 1,0\% \\
Norte de Santander & 24,3 & 0,7\% & 10,0 & 1,0\% \\
Boyacá & 47,5 & 0,7\% & 22,6 & 1,0\% \\
Atlántico & 37,3 & 0,3\% & 16,0 & 0,8\% \\
Huila & 31,6 & 0,3\% & 14,2 & 0,4\% \\
Córdoba & 22,0 & -0,6\% & 10,7 & 0,2\% \\
Magdalena & 23,9 & -0,6\% & 10,3 & 0,2\% \\
Cesar & 32,6 & -1,0\% & 13,8 & -0,2\% \\
Cundinamarca & 40,4 & -0,4\% & 17,4 & -0,2\% \\
La Guajira & 24,8 & -0,8\% & 11,1 & -0,3\% \\
\bottomrule
\end{tabular}
\caption*{\footnotesize Nota: PIB por trabajador en millones de pesos constantes de 2015; PIB por hora en miles de pesos constantes de 2015. Fuente: cálculos propios con DANE y GEIH.}
\end{table}